% Part: cheat-sheets
% Chapter: free-logic

\documentclass[../../../../include/open-logic-chapter]{subfiles}

\begin{document}

\olchapter{chs}{frl}{Free Logic}

\section{Semantics}

\subsection{Negative Free Logic}

\begin{defn}[!!^{structure}s]
In the negative system of free logic, $\Log NFL$,
\article{structure} \emph{!!{structure}}~$\Struct M$ for a language
$\Lang L_{\Log{FL}}$ of free logic consists of:
\begin{enumerate}
\item \emph{Domain:} a set, $\Domain M$,
\item \emph{Interpretation of !!{constant}s:} for any !!{constant}~$c$ of
  $\Lang L_{\Log{FL}}$, its interpretation $\Assign{c}{M}$ is !!a{element} of 
  $\Domain M$ or undefined,
\item \emph{Interpretation of !!{predicate}s:} for each $n$-place
  !!{predicate}~$R$ of $\Lang L_{\Log{FL}}$ (other than $\eq$), an $n$-place
  relation $\Assign{R}{M} \subseteq \Domain{M}^n$
\tagitem{fnTerms}{\item \emph{Interpretation of !!{function}s:} for each $n$-place
  !!{function}~$f$ of $\Lang L_{\Log{FL}}$, an $n$-place function $\Assign{f}{M}
  \colon \Domain{M}^n \to \Domain{M}$}{}
\end{enumerate}
The system of negative free logic is \emph{inclusive} if $\Domain M$ is
allowed to be empty. There is only one such !!{structure}, for interpretations
can only be defined one way when the domain is empty. We call it the 
\emph{empty !!{structure}}.
\end{defn}

\begin{defn}[Variable Assignment]
A \emph{variable assignment}~$s$ for a !!{structure}~$\Struct{M}$ is a
(partial) function such that for any !!{variable}~$x$ of
  $\Lang L_{\Log{FL}}$, $s(x) \in \Domain M$ if defined. 
\end{defn}

\begin{defn}[!!^{value} of Terms]
If $t$ is a term of the language~$\Lang L_{\Log{FL}}$, $\Struct M$ is a
!!{structure} for~$\Lang L_{\Log{FL}}$, and $s$ is a !!{variable} assignment
for~$\Struct M$, the \emph{!!{value}}~$\Value{t}{M}[s]$ is defined as
follows:
\begin{enumerate}
\item \indcase{t}{c}{$\Value{\indfrm}{M}[s] = \Assign{\indcomplex}{M}$, if defined.}
\item \indcase{t}{x}{$\Value{\indfrm}{M}[s] = s(\indcomplex)$, if defined.}
\tagitem{fnTerms}{\indcase{t}{\Atom{f}{t_1, \ldots, t_n}}{
\[
\Value{\indfrm}{M}[s] = \Assign{f}{M}(\Value{t_1}{M}[s], \ldots,
\Value{t_n}{M}[s])\textrm{, if defined}.
\]}
}{}
\end{enumerate}
\end{defn}

\begin{defn}[$x$-Variant]
If $s$ is a !!{variable} assignment for a !!{structure}~$\Struct M$, we
call !!a{variable} assignment $s'$ an \emph{$x$-variant} of $s$ iff:
\begin{enumerate}
\item it differs from $s$ at most in what it assigns for $s$, that is 
either $\Value{y}{M}[s'] = \Value{y}{M}[s]$, or both are undefined, for
any variable $y$ other than $x$.
\item it is defined for $x$, that is, $\Value{y}{M}[s'] \in \Domain M$.
We write $\varAssign{s'}{s}{x}$ when $s'$ is an $x$-variant of $s$.
\end{enumerate}
\end{defn}

\begin{defn}[Satisfaction]
\ollabel{defn:satisfaction}
Satisfaction of a !!{formula}~$!A$ in a !!{structure}~$\Struct M$
relative to a !!{variable} assignment~$s$, in symbols:
$\Sat{M}{!A}[s]$, is defined recursively as follows. (We write
$\Sat/{M}{!A}[s]$ to mean ``not $\Sat{M}{!A}[s]$.'')
\begin{enumerate}
\tagitem{prvFalse}{%
  \indcase{!A}{\lfalse}{$\Sat/{M}{\indfrm}[s]$.}}{}

\tagitem{prvTrue}{%
  \indcase{!A}{\ltrue}{$\Sat{M}{\indfrm}[s]$.}}{}

\item \indcase{!A}{\Atom{R}{t_1, \dots, t_n}}{$\Sat{M}{\indfrm}[s]$
  iff $\Value{t_i}{M}[s]$ is defined for any $1\leq i \leq n$ and 
  $\langle \Value{t_1}{M}[s], \dots, \Value{t_n}{M}[s] \rangle \in
  \Assign{R}{M}$.}

\item \indcase{!A}{\Atom{\lfrexists}{t}}{$\Sat{M}{\indfrm}[s]$ iff 
  $\Value{t}{M}[s]$ is defined.}

\item \indcase{!A}{\eq[t_1][t_2]}{$\Sat{M}{\indfrm}[s]$ iff 
  $\Value{t_1}{M}[s], \Value{t_2}{M}[s]$ are both defined and
  $\Value{t_1}{M}[s] = \Value{t_2}{M}[s]$.}

\tagitem{prvNot}{%
  \indcase{!A}{\lnot !B}{$\Sat{M}{\indfrm}[s]$ iff
    $\Sat/{M}{!B}[s]$.}}{}

\tagitem{prvAnd}{%
  \indcase{!A}{(!B \land !C)}{$\Sat{M}{\indfrm}[s]$ iff $\Sat{M}{!B}[s]$
    and $\Sat{M}{!C}[s]$.}}{}

\tagitem{prvOr}{%
  \indcase{!A}{(!B \lor !C)}{$\Sat{M}{\indfrm}[s]$ iff
    $\Sat{M}{!A}[s]$ or $\Sat{M}{!B}[s]$ (or both).}}{}

\tagitem{prvIf}{%
  \indcase{!A}{(!B \lif !C)}{$\Sat{M}{\indfrm}[s]$ iff $\Sat/{M}{!B}[s]$
    or $\Sat{M}{!C}[s]$ (or both).}}{}

\tagitem{prvIff}{%
  \indcase{!A}{(!B \liff !C)}{$\Sat{M}{\indfrm}[s]$ iff either both
    $\Sat{M}{!B}[s]$ and $\Sat{M}{!C}[s]$, or neither $\Sat{M}{!B}[s]$
    nor $\Sat{M}{!C}[s]$.}}{}

\tagitem{prvAll}{%
  \indcase{!A}{\lforall[x][!B]}{$\Sat{M}{\indfrm}[s]$ iff for every
    $x$-variant~$s'$ of $s$, $\Sat{M}{!B}[s']$.}}{}

\tagitem{prvEx}{%
  \indcase{!A}{\lexists[x][!B]}{$\Sat{M}{\indfrm}[s]$ iff there is an
    $x$-variant~$s'$ of $s$ so that $\Sat{M}{!B}[s']$.}}{}
\end{enumerate}
\end{defn}

\begin{defn}[Truth]
Where $!A$ is a closed !!{formula}, we say that $A$ is \emph{true in}, or
\emph{satisfied by} a !!{structure} $\Struct M$, written $\Sat{M}{!A}$ iff 
$A$ is satisfied relative to any assignement in $\Struct M$:
$$\Sat{M}{!A}\textrm{ iff }\Sat{M}{!A}[s]\textrm{ for any }s\textrm{ in }\Struct M$$
\end{defn}

\begin{defn}[Validity]
  A !!{formula} $!A$ is \emph{valid}, $\Entails !A$, iff $\Sat{M}{!A}$ for every
  !!{structure}~$\Struct M$.
\end{defn}

\subsection{Positive Free Logic}

\begin{defn}[!!^{structure}s]
In the positive system of free logic, $\Log NFL$,
\article{structure} \emph{!!{structure}}~$\Struct M$ for a language
$\Lang L_{\Log{FL}}$ of free logic consists of:
\begin{enumerate}
\item \emph{(Inner) Domain:} a set, $\Domain M$,
\item \emph{Outer Domain:} a non-empty set, $\Outerdomain M$ that
	includes $\Domain M$ ($\Domain M \subseteq \Outerdomain M $),
\item \emph{Interpretation of !!{constant}s:} for any !!{constant}~$c$ of
  $\Lang L_{\Log{FL}}$, its interpretation $\Assign{c}{M}$ is !!a{element} of 
  $\Outerdomain M$,
\item \emph{Interpretation of !!{predicate}s:} for each $n$-place
  !!{predicate}~$R$ of $\Lang L_{\Log{FL}}$ (other than $\eq$), an $n$-place
  relation $\Assign{R}{M} \subseteq \Outerdomain M^n$
\tagitem{fnTerms}{\item \emph{Interpretation of !!{function}s:} for each $n$-place
  !!{function}~$f$ of $\Lang L_{\Log{FL}}$, an $n$-place function $\Assign{f}{M}
  \colon \Outerdomain M^n \to \Outerdomain M$.}{}
\end{enumerate}
The system of positive free logic is \emph{inclusive} if $\Domain M$ is
allowed to be empty. There are several such !!{structure}s, as they may
still differ on the !!{element}s of $\Outerdomain M$ outside of $M$ and the 
interpretations of constants and predicates. We call a !!{structure} with
empty $\Domain M$ an \emph{empty structure}. 
\end{defn}

\begin{defn}[Variable Assignment]
A \emph{variable assignment}~$s$ for a !!{structure}~$\Struct{M}$ is a
  function such that for any !!{variable}~$x$ of
  $\Lang L_{\Log{FL}}$, $s(x) \in \Outerdomain M$. 
\end{defn}

\begin{defn}[!!^{value} of Terms]
If $t$ is a term of the language~$\Lang L_{\Log{FL}}$, $\Struct M$ is a
!!{structure} for~$\Lang L_{\Log{FL}}$, and $s$ is a !!{variable} assignment
for~$\Struct M$, the \emph{!!{value}}~$\Value{t}{M}[s]$ is defined as
follows:
\begin{enumerate}
\item \indcase{t}{c}{$\Value{\indfrm}{M}[s] = \Assign{\indcomplex}{M}$.}
\item \indcase{t}{x}{$\Value{\indfrm}{M}[s] = s(\indcomplex)$.}
\tagitem{fnTerms}{\indcase{t}{\Atom{f}{t_1, \ldots, t_n}}{
\[
\Value{\indfrm}{M}[s] = \Assign{f}{M}(\Value{t_1}{M}[s], \ldots,
\Value{t_n}{M}[s])\textrm{, if defined}.
\]}
}{}
\end{enumerate}
\end{defn}

\begin{defn}[$x$-Variant]
If $s$ is a !!{variable} assignment for a !!{structure}~$\Struct M$, we
call !!a{variable} assignment $s'$ an \emph{$x$-variant} of $s$ iff:
\begin{enumerate}
\item it differs from $s$ at most in what it assigns for $s$, that is 
 $\Value{y}{M}[s'] = \Value{y}{M}[s]$, for any variable $y$ other 
 than $x$.
\item it is defined for $x$, that is, $\Value{y}{M}[s'] \in \Domain M$.
We write $\varAssign{s'}{s}{x}$ when $s'$ is an $x$-variant of $s$.
\end{enumerate}
\end{defn}


\begin{defn}[Satisfaction]
\ollabel{defn:satisfaction}
Satisfaction of a !!{formula}~$!A$ in a !!{structure}~$\Struct M$
relative to a !!{variable} assignment~$s$, in symbols:
$\Sat{M}{!A}[s]$, is defined recursively as follows. (We write
$\Sat/{M}{!A}[s]$ to mean ``not $\Sat{M}{!A}[s]$.'')
\begin{enumerate}
\tagitem{prvFalse}{%
  \indcase{!A}{\lfalse}{$\Sat/{M}{\indfrm}[s]$.}}{}

\tagitem{prvTrue}{%
  \indcase{!A}{\ltrue}{$\Sat{M}{\indfrm}[s]$.}}{}

\item \indcase{!A}{\Atom{R}{t_1, \dots, t_n}}{$\Sat{M}{\indfrm}[s]$
  iff $\langle \Value{t_1}{M}[s], \dots, \Value{t_n}{M}[s] \rangle \in
  \Assign{R}{M}$.}

\item \indcase{!A}{\Atom{\lfrexists}{t}}{$\Sat{M}{\indfrm}[s]$ iff 
  $\Value{t}{M}[s] \in \Domain M$.}

\item \indcase{!A}{\eq[t_1][t_2]}{$\Sat{M}{\indfrm}[s]$ iff 
  $\Value{t_1}{M}[s] = \Value{t_2}{M}[s]$.}

\tagitem{prvNot}{%
  \indcase{!A}{\lnot !B}{$\Sat{M}{\indfrm}[s]$ iff
    $\Sat/{M}{!B}[s]$.}}{}

\tagitem{prvAnd}{%
  \indcase{!A}{(!B \land !C)}{$\Sat{M}{\indfrm}[s]$ iff $\Sat{M}{!B}[s]$
    and $\Sat{M}{!C}[s]$.}}{}

\tagitem{prvOr}{%
  \indcase{!A}{(!B \lor !C)}{$\Sat{M}{\indfrm}[s]$ iff
    $\Sat{M}{!A}[s]$ or $\Sat{M}{!B}[s]$ (or both).}}{}

\tagitem{prvIf}{%
  \indcase{!A}{(!B \lif !C)}{$\Sat{M}{\indfrm}[s]$ iff $\Sat/{M}{!B}[s]$
    or $\Sat{M}{!C}[s]$ (or both).}}{}

\tagitem{prvIff}{%
  \indcase{!A}{(!B \liff !C)}{$\Sat{M}{\indfrm}[s]$ iff either both
    $\Sat{M}{!B}[s]$ and $\Sat{M}{!C}[s]$, or neither $\Sat{M}{!B}[s]$
    nor $\Sat{M}{!C}[s]$.}}{}

\tagitem{prvAll}{%
  \indcase{!A}{\lforall[x][!B]}{$\Sat{M}{\indfrm}[s]$ iff for every
    $x$-variant~$s'$ of $s$, $\Sat{M}{!B}[s']$.}}{}

\tagitem{prvEx}{%
  \indcase{!A}{\lexists[x][!B]}{$\Sat{M}{\indfrm}[s]$ iff there is an
    $x$-variant~$s'$ of $s$ so that $\Sat{M}{!B}[s']$.}}{}
\end{enumerate}
\end{defn}

Truth and validity are defined as usual.

\begin{defn}[Truth]
Where $!A$ is a closed !!{formula}, we say that $A$ is \emph{true in}, or
\emph{satisfied by} a !!{structure} $\Struct M$, written $\Sat{M}{!A}$ iff 
$A$ is satisfied relative to any assignement in $\Struct M$:
$$\Sat{M}{!A}\textrm{ iff }\Sat{M}{!A}[s]\textrm{ for any }s\textrm{ in }\Struct M$$
\end{defn}

\begin{defn}[Validity]
A !!{formula} $!A$ is \emph{valid}, $\Entails !A$, iff $\Sat{M}{!A}$ for every
!!{structure}~$\Struct M$.
\end{defn}


\section{Tableaux}

\subsection{Common Rules}

The rules below are common to all systems of free logic. 

\paragraph{Rules for $\lforall$}

\begin{defish}
\AxiomC{\sFmla{\True}{\lforall[x][!A(x)]}}
\RightLabel{\TRule{\True}{\lforall}}
\UnaryInfC{$\sFmla{\False}{\lfrexists t} \quad \mid \quad \sFmla{\True}{!A(t)}$}
\DisplayProof
\hfill
\AxiomC{\sFmla{\False}{\lforall[x][!A(x)]}}
\RightLabel{\TRule{\False}{\lforall}}
\UnaryInfC{\sFmla{\True}{\lfrexists a}}
\noLine
\UnaryInfC{\sFmla{\False}{!A(a)}}
\DisplayProof
\end{defish}

In \TRule{\True}{\lforall}, $t$ is \iftag{fnTerms}{a closed term (i.e., one without
variables)}{!!a{constant}}. In \TRule{\False}{\lforall}, $a$~is !!a{constant} 
which must not occur anywhere in the branch above \TRule{\False}{\lforall}
rule. We call $a$ the \emph{eigenvariable} of the \TRule{\False}{\forall}
inference.

\paragraph{Rules for $\lexists$}

\begin{defish}
\AxiomC{\sFmla{\True}{\lexists[x][!A(x)]}}
\RightLabel{\TRule{\True}{\lexists}}
\UnaryInfC{\sFmla{\True}{\lfrexists a}}
\noLine
\UnaryInfC{\sFmla{\True}{!A(a)}}
\DisplayProof
\hfill
\AxiomC{\sFmla{\False}{\lexists[x][!A(x)]}}
\RightLabel{\TRule{\False}{\lexists}}
\UnaryInfC{$\sFmla{\False}{\lfrexists t} \quad \mid \quad \sFmla{\False}{!A(t)}$}
\DisplayProof
\end{defish}

Again, $t$~is $t$ is \iftag{fnTerms}{a closed term}{!!a{constant}}, and $a$~is 
!!a{constant} which does not occur in the branch above 
the~\TRule{\False}{\lexists} rule. We call
$a$ the \emph{eigenvariable} of the \TRule{\False}{\lexists} inference.

The condition that an eigenvariable not occur in the branch above 
the \TRule{\False}{\lforall} or \TRule{\True}{\lexists} inference is called the
\emph{eigenvariable condition}.

\paragraph{Rule for $\eq$}

The rule for $\eq$ is ($t$, $t_1$, and $t_2$ are closed terms):

\begin{defish}
\AxiomC{\sFmla{\True}{\eq[t_1][t_2]}}
\noLine
\UnaryInfC{\sFmla{\True}{!A(t_1)}}
\RightLabel{$\TRule{\True}{\eq}$}
\UnaryInfC{\sFmla{\True}{!A(t_2)}}
\DisplayProof
\hfill
\AxiomC{\sFmla{\True}{\eq[t_1][t_2]}}
\noLine
\UnaryInfC{\sFmla{\False}{!A(t_1)}}
\RightLabel{$\TRule{\False}{\eq}$}
\UnaryInfC{\sFmla{\False}{!A(t_2)}}
\DisplayProof
\end{defish}


\subsection{Negative Free Logic}

Negative free logic is obtained by supplementing the common rules with
the following rule of \emph{existential committment}. Where $R$ is some predicate,  $t$ some term, $a$ !!a{constant}, and $\ldots a \ldots$ a sequence of terms in which  $a$ appears (so $R\ldots a \ldots$ may be $Ra$, $Rbac$ or $Racacbd$, for instance):

\begin{defish}
\AxiomC{\sFmla{\True}{\Atom{R}{\ldots,a,\ldots}}}
\RightLabel{\TRule{}{\lfrexists_{\Log{NFL}}}}
\UnaryInfC{$\sFmla{\True}{\lfrexists a}$}
\DisplayProof
\hfill
\AxiomC{\sFmla{\True}{\eq[a][t]}}
\RightLabel{\TRule{}{\lfrexists_{\Log{NFL}}}}
\UnaryInfC{$\sFmla{\True}{\lfrexists a}$}
\DisplayProof
\end{defish}

\subsection{Positive Free Logic}

Positive free logic is obtained by supplementing by supplementing the 
common rules with the following rule of \emph{unrestricted self-identity}, 
where $t$ is any term:

\begin{defish}
\AxiomC{}
\RightLabel{$\eq_{\Log{PFL}}$}
\UnaryInfC{\sFmla{\True}{\eq[t][t]}}
\DisplayProof
\end{defish}


\OLEndChapterHook

\end{document}
